\subsection{Gemini 3.1 Pro mit strukturiertem Prompt}
Die im ersten Experiment (\ref{tab:Gemini_3_1_Basis_Prompt}) identifizierten Defizite bilden die Grundlage für das zweite Experiment (\ref{tab:Gemini 3.1v2}), in dem durch einen strukturierten Prompt zusätzliche Anforderungen und Kontextinformationen bereitgestellt werden.

Die Ergebnisse des zweiten Experiments (\ref{tab:Gemini 3.1v2}) unter Verwendung eines strukturierten Prompts sind in \vref{tab:Gemini 3.1v2} dargestellt. Im Vergleich zum ersten Experiment zeigt das Modell keine Leistungssteigerung. Zwar wurden grundlegende Komponenten wie eine Benutzeroberfläche, Formulare, Authentifizierung sowie teilweise Modell-Validierungen und Geschäftslogik generiert, jedoch blieben wesentliche funktionale Anforderungen unvollständig. Insbesondere wurden zentrale Geschäftsobjekte und deren vollständige CRUD-Funktionalitäten sowie erweiterte Domänenfunktionen nicht oder nur teilweise implementiert. Darüber hinaus fehlen eine vollständige Benutzer- und Rechteverwaltung, automatisierte Unit-Tests sowie End-to-End-Workflows. Auch die Frontend-Backend-Integration wurde lediglich in grundlegenden Ansätzen umgesetzt. Insgesamt erreicht das generierte System vier von 18 möglichen Bewertungspunkten und erfüllt damit 22,2\% der definierten Evaluationskriterien. 

\subsection{Gemini 3.1 Pro mit strukturiertem Prompt und Rollendefinition}
Im dritten Experiment (\ref{tab:Gemini 3.1v3}) wurde der strukturierte Prompt um eine explizite Rollendefinition erweitert, bei der das Modell die Rolle eines erfahrenen Softwarearchitekten und Full-Stack-Entwicklers einnahm. Die zusammenfassenden Ergebnisse der Evaluation sind in (\ref{tab:Gemini 3.1v3}) dargestellt.

\newpage
\begin{table}[H]
\centering
\small
\caption{Ergebnisse von Gemini 3.1 Pro mit strukturiertem Prompt}
\label{tab:Gemini 3.1v2}
\renewcommand{\arraystretch}{1.15}

\begin{tabularx}{\textwidth}{@{}>{\RaggedRight\arraybackslash}p{5.0cm} >{\centering\arraybackslash}p{1.4cm} >{\RaggedRight\arraybackslash}X@{}}
\toprule
\textbf{Evaluationskriterium} & \textbf{Wert} & \textbf{Beobachtung} \\
\midrule
\multicolumn{3}{@{}l@{}}{\textbf{Funktionale Kriterien}} \\
\addlinespace[0.2em]
Allgemeine Funktionalität & 1 &
Grundlegende Benutzeroberfläche, Formulare und Authentifizierung vorhanden. Erweiterte Funktionen wie Filterung, Pagination sowie vollständige CRUD-Operationen fehlen. \\
Geschäftsobjekte und CRUD-Funktionalität & 0 &
\textit{Customers, Orders, Rooms, Tasks und Invoices} wurden nicht umgesetzt. Mehrere Module besitzen keine vollständige CRUD-Funktionalität. \\
Erweiterte Domänenfunktionen & 0 &
\textit{\textit{Order Items, Occupancy Items, Services, Kommentare und Room Locations}} wurden nicht oder nur unvollständig implementiert. \\
Benutzer- und Rechteverwaltung & 0 &
Benutzer- und Rollenmodelle sind nicht vorhanden, jedoch fehlt eine vollständige Rechteverwaltung und Autorisierung. \\
\midrule
\multicolumn{3}{@{}l@{}}{\textbf{Technische Kriterien}} \\
\addlinespace[0.2em]
Modell-Validierungen und Geschäftslogik & 1 &
Model-Validierungen, Geschäftsregeln und technische Qualitätssicherungsmechanismen wurden teilweise umgesetzt. \\
Unit-Tests & 0 &
Es wurden keine automatisierten Unit-Tests generiert. \\
Codequalität und Wartbarkeit & 1 &
Die technische Struktur enthält teilweise nachweisbare Mechanismen zur Qualitätssicherung oder Testbarkeit. \\
\midrule
\multicolumn{3}{@{}l@{}}{\textbf{Integrationskriterien}} \\
\addlinespace[0.2em]
Frontend-Backend-Integration & 1 &
Grundlegende Interaktionen zwischen Benutzeroberfläche und Backend sind vorhanden. Komplexe Workflows wurden jedoch nicht vollständig umgesetzt. \\
End-to-End-Workflows & 0 &
Keine vollständigen Benutzerabläufe oder automatisierten Integrationstests vorhanden. \\
\midrule
\multicolumn{3}{@{}l@{}}{\textbf{Gesamtergebnis}} \\
\addlinespace[0.2em]
Erreichte Gesamtpunktzahl & 4 / 18 &
Das generierte System erreicht 4 von 18 möglichen Bewertungspunkten und erfüllt damit etwa 22,2\% der definierten Evaluationskriterien. \\
\bottomrule
\end{tabularx}
\end{table}

